\section{Calculate the conflicting sets}\label{transpose-x-sets}
In the reference layout, input starts at hexadecimal \texttt{10010000}; output
starts at \texttt{10010440}, 1088 bytes later. With sixteen words per row, each
row start advances sixty-four bytes, or four sixteen-byte line numbers.
\tagpdfsetup{table/header-rows={1}}
\begin{longtable}{rrrr}
  \toprule
  Row number & Input row byte offset & Input direct set & Output direct set \\
  \midrule
  0          & 0                     & 0                & 4                 \\
  1          & 64                    & 4                & 0                 \\
  2          & 128                   & 0                & 4                 \\
  3          & 192                   & 4                & 0                 \\
  \bottomrule
\end{longtable}
This table concerns the first line of each row, not every line in a row. Each
four-word tile row occupies one such line when its first column is a multiple of
four. Rows zero and two compete for set zero in the input; rows one and three
compete for set four. The output's row starts use those same sets in the opposite
order. Other direct-cache positions cannot hold these colliding lines. Moving
an allocation can change the pattern; record both bases when repeating the
experiment.
